\par\noindent\textbf{(1)} 胞内质量 \(m_1\) 连接 \(v_n,v_{n-1}\)，故 \(m_1\ddot u_n=K(v_n+v_{n-1}-2u_n)\)；质量 \(m_2\) 连接 \(u_n,u_{n+1}\)，\(m_2\ddot v_n=K(u_n+u_{n+1}-2v_n)\)。原胞使用长度 \(a\) 而非两倍任意坐标。
\par\noindent\textbf{(2)} 取 \(u_n=Ue^{i(qna-\omega t)}\)、\(v_n=Ve^{i(qna-\omega t)}\)，得 \((2K-m_1\omega^2)(2K-m_2\omega^2)-4K^2\cos^2(qa/2)=0\)。整理 \(\omega_\pm^2=K(1/m_1+1/m_2)\pm K\sqrt{(1/m_1+1/m_2)^2-4\sin^2(qa/2)/(m_1m_2)}\)。
\par\noindent\textbf{(3)} 当 \(qa\to0\)，低支 \(\omega_-^2\sim Ka^2q^2/[2(m_1+m_2)]\)，所以 \(c_s=a\sqrt{K/[2(m_1+m_2)]}\)。高支 \(\omega_+(0)^2=2K(1/m_1+1/m_2)\)，两质量反相运动，产生非零光学频隙。
\par\noindent\textbf{(4)} 在 \(q=\pi/a\) 时 \(\cos(qa/2)=0\)，动态矩阵对角化，两个频率平方分别为 \(2K/m_1\)、\(2K/m_2\)。若 \(m_1\ne m_2\)，两支之间存在频隙；若 \(m_1=m_2\)，区边发生折叠简并而不产生真正的质量差禁带。
